\section{Conclusion} \label{sec:conclusion}
In this work, we argue that diffusion models exhibit distinct behaviors in supervision and extrapolation regions--a perspective we call \emph{selective underfitting}, which offers a novel viewpoint on their generalization and generative performance. Theoretically, however, the mechanism enabling extrapolation remains unclear; future work may further analyze \emph{how the score in the extrapolation region is shaped}, building on our findings about the freedom of extrapolation. Practically, while we do not propose a new training method that optimally balances supervision and extrapolation, our results suggest that \emph{many well-balanced training algorithms can be developed using our decomposed analysis framework}. Ultimately, we expect these insights to extend to other generative models, contributing to a deeper understanding and the development of more powerful generative systems.


\section*{Acknowledgement}
VS was supported by the National Science Foundation under Grant No. 2211259, by the Singapore DSTA under DST00OECI20300823 (New Representations for Vision, 3D Self-Supervised Learning for Label-Efficient Vision), by the Intelligence Advanced Research Projects Activity (IARPA) via Department of Interior/ Interior Business Center (DOI/IBC) under 140D0423C0075, by the Amazon Science Hub, and by the MIT-Google Program for Computing Innovation. SC was supported by NSF SLES award IIS-2331831 and the Harvard Dean's Competitive Fund for Promising Scholarship.